\chapter{矩阵：运算、求逆与秩}\label{ch:matrix}
\chaptermap{运算与伴随}{求逆四法}{秩与分解}{方程与分块}
\begin{examtip}
\textbf{考研核心：}矩阵运算、逆矩阵、初等矩阵、秩、矩阵方程。每次运算先写维数，再写乘法顺序。
\textbf{竞赛视角：}用多项式关系降次，用秩分解压缩维数，用分块消元统一求逆与解方程。经典伴随记作$\adj A$，不与共轭转置混用。
\end{examtip}
\section{运算规则与经典伴随}
\subsection{乘法次序是条件的一部分}
\begin{definition}{乘积、转置与伴随}
矩阵相等要求形状相同且对应元素相同。同形矩阵按元素相加、数乘：$(A+B)_{ij}=a_{ij}+b_{ij}$，$(cA)_{ij}=ca_{ij}$。若$A\in\F^{m\times n},B\in\F^{n\times p}$，乘积为$m\times p$矩阵，$(AB)_{ij}=\sum_{k=1}^n a_{ik}b_{kj}$。转置由$(A\T)_{ij}=a_{ji}$定义。

$I_n$的对角元为1、其余为0。对$n$阶$A$，若存在同阶$B$使$AB=BA=I_n$，称$A$可逆，$B$称其逆矩阵，记$A\inv$；其唯一性由$B=B(AC)=(BA)C=C$得到。方阵的迹定义为$\tr A=\sum_i a_{ii}$。实方阵满足$A\T=A$称对称，满足$A\T=-A$称反对称；实方阵$Q$满足$Q\T Q=I$称正交矩阵，此时$Q\inv=Q\T$，也有$QQ\T=I$。

对方阵，经典伴随为$\adj A=(C_{ji})$。三类初等行变换是交换两行、一行乘非零标量、某行加上另一行的标量倍；初等列变换相应定义。对$I_n$作一次初等行变换所得矩阵称初等矩阵，左乘它实施相同的行变换；右乘相应初等矩阵实施列变换。
\end{definition}
\begin{theorem}{代数规则与伴随恒等式}
乘法满足结合律与分配律，通常不满足交换律。并有
\begin{equation}(AB)\T=B\T A\T,\quad A\adj A=\adj A\,A=(\det A)I.\label{eq:adj}\end{equation}
可逆时$(AB)\inv=B\inv A\inv$，$(A\T)\inv=(A\inv)\T$。$n$阶矩阵满足$\det(\adj A)=(\det A)^{n-1}$；$n=1$时伴随按$[1]$约定。
\end{theorem}
\begin{proof}
\label{der:2-1}
先说明迹交换公式：若$A\in\F^{m\times n},B\in\F^{n\times m}$，则
\[\tr(AB)=\sum_{i=1}^m\sum_{j=1}^n a_{ij}b_{ji}=\sum_{j=1}^n\sum_{i=1}^m b_{ji}a_{ij}=\tr(BA).\]
特别$\tr(P\inv AP)=\tr(APP\inv)=\tr A$，所以相似保持迹；一般合同不具备这一保证。

结合律逐项为
\[((AB)C)_{ij}=\sum_\ell\sum_k a_{ik}b_{k\ell}c_{\ell j}=\sum_k a_{ik}\sum_\ell b_{k\ell}c_{\ell j}=(A(BC))_{ij}.\]
分配律由有限和的分配律逐项得到。转置公式的$(i,j)$元为
\[((AB)\T)_{ij}=\sum_k a_{jk}b_{ki}=\sum_k(B\T)_{ik}(A\T)_{kj}=(B\T A\T)_{ij}.\]
伴随恒等式左乘积的$(i,j)$元为$\sum_k a_{ik}C_{jk}=\delta_{ij}\det A$，右乘积按列的同列、异列展开同样得到$\delta_{ij}\det A$，因此两边均为$(\det A)I$。

若$A,B$可逆，直接计算
\[(AB)(B\inv A\inv)=A(BB\inv)A\inv=I,\quad
(B\inv A\inv)(AB)=B\inv(A\inv A)B=I,\]
故候选矩阵就是乘积的逆。将$AA\inv=I,A\inv A=I$分别转置，得到$(A\inv)\T A\T=A\T(A\inv)\T=I$，证明转置求逆公式。

可逆时$\adj A=(\det A)A\inv$，于是
\[\det(\adj A)=(\det A)^n\det A\inv=(\det A)^{n-1}.\]
若$A$奇异，考察$A+tI$。$\det(A+tI)$是首一$n$次多项式，只在有限个$t$上为零；两边关于$t$均为多项式，在其余无限多个$t$上相等，所以作为多项式恒等，取$t=0$得到奇异情形。$n=1$时直接由$\adj[a]=[1]$核验，右边按零次幂恒为1解释。

\end{proof}
\begin{method}{先整理等式再计算}
适用：复合矩阵式。步骤：核对维数；按括号逆序展开逆与转置；只对可交换的同一矩阵多项式合并。易错：不能任意约去奇异矩阵，$AB=0$也不能推出某个因子为零。
\end{method}
\worked{不交换的最小例子}{基础例题}{\star}{2-1}
{取$A=\begin{pmatrix}0&1\\0&0\end{pmatrix}$，$B=\begin{pmatrix}0&0\\1&0\end{pmatrix}$，求$AB,BA$。}
{按“行乘列”计算。}
{逐项行乘列：
\[AB=\begin{pmatrix}0\cdot0+1\cdot1&0\cdot0+1\cdot0\\0&0\end{pmatrix}=\begin{pmatrix}1&0\\0&0\end{pmatrix},\quad
BA=\begin{pmatrix}0&0\\1\cdot0+0\cdot0&1\cdot1+0\cdot0\end{pmatrix}=\begin{pmatrix}0&0\\0&1\end{pmatrix}.\]
两矩阵的$(1,1)$元不同，故不交换。另算$A^2=B^2=O$，但$AB\ne O$，因此幂零不意味着与别的矩阵相乘为零。}
{一个二阶反例就能否定无条件的交换与消去。}
{展开$(A+B)^2$应保留$AB+BA$，不能默认是$2AB$。}
\worked{逆与转置的顺序}{考研风格}{\star\star}{2-2}
{可逆$A,B$满足$X=(A\T B)\inv$，写出$X\T$。}
{每遇一次逆或转置，整体因子顺序反转。}
{\textbf{方法一：}先展开逆，再转置：$X=B\inv(A\T)\inv=B\inv(A\inv)\T$，于是
\[X\T=\bigl((A\inv)\T\bigr)\T(B\inv)\T=A\inv(B\inv)\T.\]
\textbf{方法二：}先把整体逆与转置交换，再转置内部乘积：
\[X\T=\bigl((A\T B)\T\bigr)\inv=(B\T A)\inv=A\inv(B\T)\inv=A\inv(B\inv)\T.\]
每一步都保留因子次序，且可逆性由$A,B$可逆保证。}
{两条路线得到相同的乘法次序，可用于自检。}
{相同字母位置不同也不能跨越另一个矩阵约去。}
\worked{伴随需要转置}{基础例题}{\star}{2-3}
{求$A=\begin{pmatrix}1&2\\3&6\end{pmatrix}$的伴随并核验式\eqref{eq:adj}。}
{二阶伴随是对角交换、副对角变号。}
{二阶代数余子式为$C_{11}=6,C_{12}=-3,C_{21}=-2,C_{22}=1$，转置后
\[\adj A=\begin{pmatrix}6&-2\\-3&1\end{pmatrix}.\]
核验$\det A=1\cdot6-2\cdot3=0$，并逐项相乘：
\[A\adj A=\begin{pmatrix}6-6&-2+2\\18-18&-6+6\end{pmatrix}=O=(\det A)I.\]
伴随第二行是第一行的$-\frac{1}{2}$倍且不为零，故秩1；原矩阵虽不可逆，伴随仍有定义。}
{奇异矩阵仍有伴随，只有求逆时才需除以行列式。}
{代数余子式矩阵未转置时不是经典伴随。}
\worked{幂等关系的逆}{考研风格}{\star\star}{2-4}
{若$P^2=P$，证明$I+P$可逆并求逆。}
{寻找$I+cP$使乘积退化为$I$。}
{\textbf{方法一：}待定逆$I+cP$，利用$P^2=P$得$(I+P)(I+cP)=I+(1+2c)P$。令$1+2c=0$，取$c=-\frac{1}{2}$，左右乘均得到$I$，故逆为$I-\frac{P}{2}$。
\textbf{方法二：}任意$x=Px+(I-P)x$，前一项在$\im P$，后一项在$\ker P$；交中向量$z$同时满足$Pz=z$与$Pz=0$，故只有零。在前空间$I+P$乘2，在后空间乘1，逆应分别乘$\frac{1}{2},1$，即$\frac12P+(I-P)=I-\frac{P}{2}$。}
{系数待定法快，空间分解法解释为何只需两个系数。}
{一般$(I+P)\inv$不等于$I-P$。}
\begin{insight}{逆矩阵是把变形倒着做}
矩阵的每一列是标准基向量变换后的去向。对$S=\begin{pmatrix}1&1\\0&1\end{pmatrix}$，竖直单位边变成$(1,1)\T$，水平边不动，单位正方形被剪成平行四边形。逆变换把每个点的横坐标减去其纵坐标，因此$S\inv=\begin{pmatrix}1&-1\\0&1\end{pmatrix}$。
\end{insight}
\begin{figure}[H]\centering
\begin{tikzpicture}[scale=1.35]
\begin{scope}
\draw[laaxis](-.2,0)--(2.4,0);\draw[laaxis](0,-.2)--(0,1.6);
\filldraw[fill=Teal!12,draw=Teal](0,0)rectangle(1,1);
\draw[lavec,Ink](0,0)--(1,0)node[midway,below]{$e_1$};\draw[lavec,Gold](0,0)--(0,1)node[left]{$e_2$};
\node[lanote]at(.7,-.5){原坐标};
\end{scope}
\draw[-{Stealth},Teal,thick](2.5,.75)--(3.8,.75)node[midway,above]{$S$};
\draw[-{Stealth},Gold,thick](3.8,.25)--(2.5,.25)node[midway,below]{$S\inv$};
\begin{scope}[xshift=4.3cm]
\draw[laaxis](-.2,0)--(2.4,0);\draw[laaxis](0,-.2)--(0,1.6);
\filldraw[fill=Teal!12,draw=Teal](0,0)--(1,0)--(2,1)--(1,1)--cycle;
\draw[lavec,Ink](0,0)--(1,0)node[midway,below]{$Se_1$};\draw[lavec,Gold](0,0)--(1,1)node[above]{$Se_2$};
\node[lanote]at(.9,-.5){剪切后的坐标};
\end{scope}
\end{tikzpicture}
\caption{先作$S$再作$S\inv$，每个向量回到原处，而不只是面积恢复。}\label{fig:matrix-inverse}
\end{figure}
\begin{clevertrick}{用幂零看穿剪切矩阵}
写$S=I+N$，其中$N=\begin{pmatrix}0&1\\0&0\end{pmatrix}$，直接相乘得$N^2=0$。当$m\ge0$时，$I,N$可交换，二项式展开中的二次及以上$N$项全为零，故
\[S^m=\sum_{j=0}^m\binom mjN^j=I+mN=\begin{pmatrix}1&m\\0&1\end{pmatrix}.\]
又$(I+N)(I-N)=I-N^2=I$，反序相乘也为$I$，所以$S\inv=I-N$。若$m=-k<0$，则$S^m=(I-N)^k=I-kN=I+mN$，幂公式仍成立。例如
\[S^{100}(0,1)\T=(0,1)\T+100N(0,1)\T=(0,1)\T+100(1,0)\T=(100,1)\T.\]
一般块矩阵$S=I+N$若只有一个非零副对角块，先核验$N^2=0$，便可逐字使用同一推导。
\textbf{易错提示：}面积不变不表示矩阵等于单位阵；图中剪切的行列式为1，但大部分向量都改变了方向。
\end{clevertrick}


\section{逆矩阵的四条路线}
\subsection{伴随、消元、分块与多项式}
\begin{theorem}{可逆的等价刻画}
对$n$阶$A$，下列等价：$A$可逆；$\det A\ne0$；$\rank A=n$；$A\vx=0$只有零解；$A$行等价于$I$；$A$为有限个初等矩阵之积。此时
\begin{equation}A\inv=\dfrac{\adj A}{\det A},\qquad [A\mid I]\longrightarrow[I\mid A\inv].\end{equation}
\end{theorem}
\begin{proof}
\label{der:2-2}
若$A$可逆，$A\vx=0$左乘$A\inv$给$\vx=0$。反过来，对$A$作初等行变换得到阶梯形$U=EA$，其中$E$是可逆初等矩阵的乘积。每种行操作都有逆操作：换行再换回，倍乘$c\ne0$的逆是倍乘$\frac{1}{c}$，倍加$c$的逆是倍加$-c$，所以$A\vx=0$与$U\vx=0$同解。

若方阵的主元少于$n$，至少有一个自由变量；令一个自由变量为1、其余为0，回代得到非零齐次解，与假设矛盾。因此每列均为主元列，每行也非零，继续把主元归一并消去上方元素，存在初等矩阵乘积$F$使$FA=I$。于是$A=F\inv$，其逆为$F$。这同时证明行等价于$I$、可逆、满秩、无非零齐次解彼此等价，也证明$A$能写成有限个初等矩阵的乘积。反向，初等矩阵均可逆，其有限乘积可逆。

由乘积行列式，$A$可逆推出$\det A\det A\inv=1$，所以行列式非零；若行列式非零，伴随恒等式给
\[A\dfrac{\adj A}{\det A}=\dfrac{\adj A}{\det A}A=I,\]
故可逆。若两个逆为$B,C$，则$B=B(AC)=(BA)C=C$，逆唯一。

对增广矩阵，每个左乘初等矩阵同时作用于两块：$[A\mid I]\mapsto[FA\mid F]$。当左块为$I$时$F=A\inv$，所以右块恰为逆；这解释了同步行变换法而不只是给出算法。

\end{proof}
\begin{method}{求逆按结构选路}
二、三阶精确题可用伴随；一般数字题用增广矩阵；块三角先求对角块逆；给$A$的低次关系则先凑$I$。易错：行变换必须同步作用于右块，不可在$[A\mid I]$中随意混入列变换。
\end{method}
\worked{同一矩阵的四法求逆}{考研风格}{\star\star}{2-5}
{求$A=\begin{pmatrix}2&1\\1&1\end{pmatrix}$的逆。}
{本题既适合二阶公式，也可演示一般方法。}
{\textbf{方法一（伴随）：}$\det A=2-1=1$，伴随$\begin{pmatrix}1&-1\\-1&2\end{pmatrix}$，除以1即为逆。
\textbf{方法二（消元）：}
\[\left[\begin{array}{cc|cc}2&1&1&0\\1&1&0&1\end{array}\right]
\xrightarrow{r_1-r_2}\left[\begin{array}{cc|cc}1&0&1&-1\\1&1&0&1\end{array}\right]
\xrightarrow{r_2-r_1}\left[\begin{array}{cc|cc}1&0&1&-1\\0&1&-1&2\end{array}\right].\]
右块就是所求。
\textbf{方法三（分块）：}主块为标量2，Schur补$S=1-\frac{1}{2}=\frac{1}{2}$，其逆为2。四块分别为$\frac{1}{2}+(\frac{1}{2})\cdot1\cdot2\cdot1\cdot(\frac{1}{2})=1$、$-(\frac{1}{2})\cdot1\cdot2=-1$、$-2\cdot1\cdot(\frac{1}{2})=-1$、$2$。
\textbf{方法四（CH）：}$\chi_A(t)=(t-2)(t-1)-1=t^2-3t+1$，所以$A(3I-A)=I$，逆为$3I-A=\begin{pmatrix}1&-1\\-1&2\end{pmatrix}$。乘回得$\begin{pmatrix}2-1&-2+2\\1-1&-1+2\end{pmatrix}=I$，四法一致。}
{小阶数伴随最省；消元最通用；分块重用主块；多项式法适合抽象矩阵。}
{四种方法都应以$AA\inv=I$作最后核验。}
\worked{幂零矩阵的有限几何级数}{竞赛风格}{\star\star}{2-6}
{若$N^k=0$，求$(I-N)\inv$。}
{模仿有限等比和，不需要无穷级数收敛。}
{令$S=I+N+\cdots+N^{k-1}$，相乘后按同次幂相减：
\[(I-N)S=(I+N+\cdots+N^{k-1})-(N+N^2+\cdots+N^k)=I-N^k=I.\]
所有项是$N$的多项式，可交换，所以$S(I-N)=I$也成立，$S$为双侧逆。例如取$N=\begin{pmatrix}0&1&0\\0&0&1\\0&0&0\end{pmatrix}$，则$N^2=\begin{pmatrix}0&0&1\\0&0&0\\0&0&0\end{pmatrix},N^3=0$，故$I+N+N^2=\begin{pmatrix}1&1&1\\0&1&1\\0&0&1\end{pmatrix}$。这是有限和，不需要讨论无穷级数收敛。}
{多项式恒等式比逐项求伴随高效。}
{对一般$N$不能把无穷几何级数无条件代入。}
\worked{由关系式构造逆}{考研风格}{\star\star}{2-7}
{已知$A^2-A-2I=0$，求$A\inv$与$(I-A)\inv$。}
{从原式提取乘积等于常数单位阵。}
{\textbf{方法一：}原式给$A^2-A=2I$，即$A(A-I)=(A-I)A=2I$，故$A\inv=\frac{A-I}{2}$。另有$(I-A)A=A-A^2=-2I$，所以$(I-A)\inv=-\frac{A}{2}$。
\textbf{方法二：}设$A\inv=qA+rI$，利用$A^2=A+2I$，有$A(qA+rI)=(q+r)A+2qI$；取$q=\frac{1}{2},r=-\frac{1}{2}$使之为$I$。对$(I-A)\inv=uA+vI$，乘开并降次得到$-vA+(v-2u)I$，取$v=0,u=-\frac{1}{2}$。两次都是构造出逆后再核验，不预设未知矩阵可逆。}
{凑单位阵最直接，模多项式方法可处理更高次题。}
{$(A+I)(A-2I)=0$只推出两因子不能同时可逆。}
\worked{块上三角求逆}{考研风格}{\star\star}{2-8}
{设$A,D$可逆，求$M=\begin{pmatrix}A&B\\0&D\end{pmatrix}$的逆。}
{未知逆也设为块上三角。}
{\textbf{方法一：}设候选逆为$\begin{pmatrix}X&Y\\0&Z\end{pmatrix}$，乘积为$\begin{pmatrix}AX&AY+BZ\\0&DZ\end{pmatrix}$。等于$I$要求$X=A\inv,Z=D\inv,Y=-A\inv BD\inv$，得到
\[M\inv=\begin{pmatrix}A\inv&-A\inv BD\inv\\0&D\inv\end{pmatrix}.\]
反向相乘的右上块为$A\inv B-A\inv BD\inv D=0$，也等于$I$。
\textbf{方法二：}$M=\diag(A,D)(I+K)$，$K=\begin{pmatrix}0&A\inv B\\0&0\end{pmatrix}$满足$K^2=0$。所以$M\inv=(I-K)\diag(A\inv,D\inv)$，逐块乘开正是上式。}
{待定块法通用，分解法解释负号与次序。}
{右上角通常不能写成$-B A\inv D\inv$。}
\section{秩、初等变换与秩分解}
\subsection{子式与空间描述相互验证}
\begin{definition}{秩与等价标准形}
$\rank A$是非零子式的最大阶数；零矩阵的秩为零。各行向量中极大无关子组的向量数称行秩，各列相应称列秩。两个$m\times n$矩阵$A,B$若存在可逆$m$阶$P$及可逆$n$阶$Q$使$B=PAQ$，称矩阵等价；只用左侧初等行变换联系时称行等价。矩阵等价区别于相似与合同。
\end{definition}
\begin{theorem}{秩的基本性质}
行秩、列秩与非零子式最大阶数相等。可逆左、右乘保持秩；任意秩$r$的$m\times n$矩阵等价于$\begin{pmatrix}I_r&0\\0&0\end{pmatrix}$，称秩的等价标准形。乘积秩有上界；实矩阵的Gram矩阵与原矩阵同秩：
\[\rank(AB)\le\min(\rank A,\rank B),\qquad\rank(A\T A)=\rank A.\]
若$n\ge2$，经典伴随矩阵满足
\[\rank(\adj A)=\begin{cases}n,&\rank A=n,\\1,&\rank A=n-1,\\0,&\rank A\le n-2.\end{cases}\]
\end{theorem}
\begin{proof}
\label{der:2-3}
\textbf{秩的三种定义一致。}可逆行变换保列间关系，因$EA\vx=0\Longleftrightarrow A\vx=0$；列变换把列组作可逆线性组合，保列空间维数。对行空间反向论证也成立。取任一非零元作主元，换到左上角，归一并用行、列倍加清掉同列同行的其余元，对右下块重复，得到$PAQ=\diag(I_r,0)$的矩形版本。该标准形行、列独立方向都恰为$r$。

还须说明非零子式的最大阶数未变。对任意一次行操作，任一$k$阶新子式按被改变的行多线性展开，是旧矩阵$k$阶子式的线性组合（重复行项为零）；故若全部旧$k$阶子式为零，新子式也全零。逆行操作给反向蕴含；列操作同理。因此最大非零子式阶数不变，标准形中恰为$r$。行秩、列秩、子式秩由此相等。

\textbf{乘积与Gram秩。}$\im AB\subseteq\im A$给$\rank AB\le\rank A$；映射$A$把$\im B$送到$\im AB$，输出维数不超过输入维数，给$\rank AB\le\rank B$。实数域中，$A\vx=0$直接推出$A\T A\vx=0$；反向若$A\T A\vx=0$，则
\[0=\vx\T A\T A\vx=\sum_i(A\vx)_i^2,\]
每项非负，所以每项为零。两矩阵核相同，未知量数相同，故秩相同。

\textbf{伴随秩三类。}若$\rank A=n$，$\adj A=(\det A)A\inv$可逆，秩$n$。若$\rank A=n-1$，存在非零$(n-1)$阶子式，所以$\adj A\ne0$；又$A\adj A=0$，伴随每列在一维$\ker A$中，故其秩至多1且非零，只能为1。若$\rank A\le n-2$，全部$(n-1)$阶子式为零，所以伴随为零。这里$n\ge2$确保余子式阶数至少1。

\end{proof}
\begin{method}{秩的上下界夹逼}
先找非零子式得下界，再找相关关系得上界；参数题先求满阶行列式的零点，再在这些点检查低阶子式。不要用“行列式为零”直接判断秩降了几级。
\end{method}
\worked{参数秩的三个区间}{考研风格}{\star\star}{2-9}
{三阶矩阵对角元为$t$、非对角元为$1$，求秩。}
{由第1章公式先找到特殊参数。}
{\textbf{方法一：}行列式为$(t-1)^2(t+2)$。$t\ne1,-2$时非零，秩3；$t=1$时所有列均为$\bm1$，有非零元且列同向，秩1；$t=-2$时行列式零，秩至多2，而左上二阶子式$(-2)^2-1=3\ne0$给下界2，故秩2。
\textbf{方法二：}$A=(t-1)I+\bm1\bm1\T$。在$\Span(\bm1)$上特征值为$t+2$，在二维和零平面上为$t-1$。这两个空间直和为$\R^3$，故秩等于非零标量作用所覆盖的维数：一般为$1+2=3$，$t=1$只剩前一维，$t=-2$只剩后两维。此计数依赖已经证明的直和结构。}
{子式法不依赖可对角化；谱法适合这个特殊结构。}
{仅看行列式的二重零点，不能对一般矩阵推断秩降$2$。}
\worked{伴随矩阵的秩}{考研风格}{\star}{2-10}
{四阶$A$的秩依次为$4,3,2$时，求伴随的秩。}
{使用按$n-1$阶子式划分的三类。}
{若原秩4，则$A$可逆，$\adj A=(\det A)A\inv$也可逆，秩4。若原秩3，某个三阶子式非零使伴随非零；又$A\adj A=0$，伴随全部列在一维核中，因此秩1。若原秩2，全部三阶子式为零，伴随所有元素为零，秩0。三个答案依次是$4,1,0$，不是把原秩统一减1。}
{这里的$n-1$是结构阈值，不是“秩减一”。}
{当$n=1$时经典伴随恒为$[1]$，需单独看约定。}
\worked{实数域中的Gram矩阵}{竞赛风格}{\star\star}{2-11}
{证明实矩阵$A$满列秩当且仅当$A\T A$可逆，并说明复数域要如何修改。}
{把二次式变成平方和。}
{\textbf{方法一：}对任意实$x$，$x\T A\T Ax=\sum_i(Ax)_i^2$。所以$A\T Ax=0$推出平方和为零，进而$Ax=0$；反向由矩阵乘法直接得到。两核相同，所以$A$满列秩恰当于$A\T A$核为零，也恰当于它作为方阵可逆。
\textbf{方法二：}若列无关，对非零系数$x$有$Ax\ne0$，故Gram二次型$x\T A\T Ax>0$，正定使Gram矩阵可逆；若列相关，取非零$x$使$Ax=0$，便给Gram的非零核向量，不能可逆。复数域须改用$x^*A^*Ax=\sum_i|(Ax)_i|^2$；只取转置时$A=(1,i)\T$虽满列秩，却有$A\T A=1+i^2=0$。}
{平方和论证同时给出可逆性和正定性。}
{复数的平方和可以抵消，共轭内积才有正性。}
\worked{秩分解}{竞赛风格}{\star\star}{2-12}
{将$A=\begin{pmatrix}1&0&1\\0&1&1\\1&1&2\end{pmatrix}$写成满列秩矩阵与满行秩矩阵的乘积。}
{第三列等于前两列之和。}
{\textbf{方法一：}前两列$b_1=(1,0,1)\T,b_2=(0,1,1)\T$独立，第三列等于$b_1+b_2$，故
\[B=\begin{pmatrix}1&0\\0&1\\1&1\end{pmatrix},\quad C=\begin{pmatrix}1&0&1\\0&1&1\end{pmatrix},\quad BC=(b_1,b_2,b_1+b_2)=A.\]
$B$前二行含$I_2$，$C$前二列含$I_2$，两者秩均2。
\textbf{方法二：}对$A$作$r_3\leftarrow r_3-r_1-r_2$，得$\begin{pmatrix}1&0&1\\0&1&1\\0&0&0\end{pmatrix}$，主元列为1、2；由非主元列的系数$(1,1)\T$恢复原列表示，仍得到上述$B,C$。必须返回原列，不能把消元后的列误作$B$。}
{秩分解把大矩阵运算压到$r$维，是低秩技巧的基础。}
{采用本题的选列算法时，应取原矩阵中对应主元编号的列，不能直接搬用行变换后的列。一般秩分解中，$B$也可取原列空间的任意一组基。}
\begin{insight}{秩是在问还剩几个独立方向}
二维图形经满秩线性变换仍可有两个独立方向；秩1时整个平面压到一条过原点的直线；秩0时只剩原点。因此$\det A=0$只说二维面积消失，并未说明最后剩下一条线还是一个点。三阶及更高阶的情形也要用秩进一步分辨。
\end{insight}
\begin{figure}[H]\centering
\begin{tikzpicture}[scale=1.05]
\begin{scope}
\draw[laaxis](-.3,0)--(3,0);\draw[laaxis](0,-.3)--(0,2);
\filldraw[fill=Teal!10,draw=Teal](0,0)--(2,0)--(2.5,1.4)--(.5,1.4)--cycle;
\draw[lavec,Ink](0,0)--(2,0);\draw[lavec,Gold](0,0)--(.5,1.4);
\node[lanote]at(1.3,-.65){两个独立方向：秩2};
\end{scope}
\draw[-{Stealth},Rust,thick](3.3,.8)--(4.8,.8)node[midway,above]{$\diag(1,0)$};
\begin{scope}[xshift=5.5cm]
\draw[laaxis](-.3,0)--(3,0);\draw[laaxis](0,-.3)--(0,2);
\draw[Teal!35,line width=3pt](0,0)--(2.5,0);
\draw[lavec,Ink](0,0)--(2,0)node[midway,above]{$A\vu$};\draw[lavec,Gold](0,0)--(.5,0)node[midway,below]{$A\vv$};
\node[lanote]at(1.3,-.65){只剩一个方向：秩1};
\end{scope}
\end{tikzpicture}
\caption{秩描述像空间维数；压缩方向构成核，满足秩与零化度之和等于原维数。}\label{fig:rank-collapse}
\end{figure}
\begin{clevertrick}{低秩矩阵先写成外积}
矩阵$K=\begin{pmatrix}1&2&3\\2&4&6\end{pmatrix}$的两行成比例，可写$K=(1,2)\T(1,2,3)$。因此$Kx=(1,2)\T(x_1+2x_2+3x_3)$，所有输出均在$\Span((1,2)\T)$，且取$x=e_1$确实得到非零输出，故秩恰为1。
若$K=\vu\vv\T$为方阵，记$c=\vv\T\vu$，结合律给
\[K^2=\vu(\vv\T\vu)\vv\T=cK.\]
由$K^1=K$起归纳，$K^m=c^{m-1}K$对$m\ge1$成立；$c=0$时应分别写$K^1=K,K^m=0\ (m\ge2)$，避免在记号上误解零次幂。又$K\vu=c\vu$、$Kx=0$对$\vv\T x=0$成立，直接给出相关特征方向。求$(I+K)\inv$时试$I-\alpha K$，乘积为$I+[1-\alpha(1+c)]K$；若$1+c\ne0$，取$\alpha=\frac{1}{1+c}$即得逆的更新公式。
\textbf{易错提示：}外积中的两个向量必须非零，才能断言秩恰为1；写出一个分解$A=BC$，一般只能先得秩的上界。
\end{clevertrick}


\section{矩阵方程与分块技巧}
\subsection{左、右乘与Schur逆}
\begin{theorem}{矩阵方程与分块逆}
若$A,B$可逆，则$AXB=C$唯一解为$X=A\inv CB\inv$。若$A$可逆且$S=D-CA\inv B$可逆，则
\begin{equation}
\begin{pmatrix}A&B\\C&D\end{pmatrix}^{-1}
=\begin{pmatrix}A\inv+A\inv BS\inv CA\inv&-A\inv BS\inv\\-S\inv CA\inv&S\inv\end{pmatrix}.\label{eq:blockinverse}
\end{equation}
\end{theorem}
\begin{proof}
\label{der:2-4}
对$AXB=C$，依次左乘$A\inv$、右乘$B\inv$，得到$X=A\inv CB\inv$；代回得$A(A\inv CB\inv)B=C$，说明存在。任何解都必须经同样两步变成此式，说明唯一。

令$M=\begin{pmatrix}A&B\\C&D\end{pmatrix}$及$S=D-CA\inv B$。先作块消元，
\[LM=U,\quad L=\begin{pmatrix}I&0\\-CA\inv&I\end{pmatrix},\quad U=\begin{pmatrix}A&B\\0&S\end{pmatrix}.\]
若$A,S$都可逆，设$U\inv=\begin{pmatrix}X&Y\\0&Z\end{pmatrix}$，由$UU\inv=I$逐块得到$AX=I,SZ=I,AY+BZ=0$，因而
\[U\inv=\begin{pmatrix}A\inv&-A\inv BS\inv\\0&S\inv\end{pmatrix}.\]
由$LM=U$有$M\inv=U\inv L$。逐块乘开：
\[\begin{aligned}
(M\inv)_{11}&=A\inv+(-A\inv BS\inv)(-CA\inv),\\
(M\inv)_{12}&=-A\inv BS\inv,\\
(M\inv)_{21}&=-S\inv CA\inv,\qquad (M\inv)_{22}=S\inv.
\end{aligned}\]
这就是给出的四块逆式。两个负号产生左上角的加号，因子顺序始终来自块乘法，不能任意交换。

\end{proof}
\begin{method}{方程分列与结构消元}
$AX=C$可按$C$的列逐列求解；$XA=C$可转置处理。若出现$AX+XB$，不要误作$(A+B)X$，可设元素或进入特征基。块逆公式较长，优先现场消元以免记错次序。
\end{method}
\worked{双侧矩阵方程}{考研风格}{\star\star}{2-13}
{设$A=\diag(2,1),B=\diag(1,3)$，解$AXB=\begin{pmatrix}2&6\\1&3\end{pmatrix}$。}
{对角左乘按行缩放，右乘按列缩放。}
{\textbf{方法一：}左乘$A\inv=\diag(\frac{1}{2},1)$先将右端按行缩放，得$\begin{pmatrix}1&3\\1&3\end{pmatrix}$；右乘$B\inv=\diag(1,\frac{1}{3})$再按列缩放，得到$X=\begin{pmatrix}1&1\\1&1\end{pmatrix}$。
\textbf{方法二：}设元素为$x_{ij}$，四式依次为$2x_{11}=2,6x_{12}=6,x_{21}=1,3x_{22}=3$，全解得1。代回$AX=\begin{pmatrix}2&2\\1&1\end{pmatrix}$，再右乘$B$恢复题给右端，且四个非零系数保证唯一。}
{元素法适合对角系数，逆矩阵公式更一般。}
{右乘逆会缩放列，不能把它施加到行。}
\worked{奇异系数下的矩阵方程}{考研风格}{\star\star}{2-14}
{设$A=\begin{pmatrix}1&1\\2&2\end{pmatrix}$，解$AX=\begin{pmatrix}1&2\\2&4\end{pmatrix}$。}
{每一列是一个相容但不唯一的方程组。}
{设$X=\begin{pmatrix}x&y\\z&w\end{pmatrix}$。乘开有$AX=\begin{pmatrix}x+z&y+w\\2x+2z&2y+2w\end{pmatrix}$，所以独立条件仅$x+z=1,y+w=2$，第二行自动为第一行的两倍。令$z=s,w=t$，解得$x=1-s,y=2-t$，因此
\[X=\begin{pmatrix}1-s&2-t\\s&t\end{pmatrix},\qquad s,t\in\R.\]
任意$s,t$代回均满足四个元素等式；任何解也都能用自己的$z,w$作$s,t$，故没有遗漏。}
{奇异不等于无解；相容时每一列都可添加核向量。}
{不能使用不存在的$A\inv$。若右端第二行不为第一行的两倍，则无解。}
\worked{一般分块逆的实例}{竞赛风格}{\star\star\star}{2-15}
{求$M=\begin{pmatrix}I&I\\I&2I\end{pmatrix}$的逆，各$I$阶数相同。}
{Schur补是$I$，也可把未知向量分块。}
{\textbf{方法一：}主块$I$可逆，Schur补$2I-II^{-1}I=I$，将四块逆式代入得$M\inv=\begin{pmatrix}2I&-I\\-I&I\end{pmatrix}$。核验左上块$2I-I=I$，右上块$-I+I=0$，左下块$2I-2I=0$，右下块$-I+2I=I$。
\textbf{方法二：}解$M(u,v)\T=(a,b)\T$，即$u+v=a,u+2v=b$。相减得$v=b-a$，再得$u=a-v=2a-b$，所以逆映射的四块分别为$2I,-I,-I,I$。两方法均不要求各块是一阶数，但要求尺寸相同。}
{解方程法更不易记错，适合验证公式。}
{两个块行都必须满足原方程。}
\worked{低秩更新的逆}{竞赛风格}{\star\star\star}{2-16}
{证明当$1+\vv\T\vu\ne0$时，$(I+\vu\vv\T)\inv=I-\frac{\vu\vv\T}{1+\vv\T\vu}$。}
{平方后仍是原秩一矩阵的倍数。}
{\textbf{方法一：}记$K=uv\T,c=v\T u$，则$K^2=u(v\T u)v\T=cK$。因此
\[(I+K)\left(I-\dfrac{K}{1+c}\right)=I+\left(1-\dfrac1{1+c}-\dfrac c{1+c}\right)K=I.\]
因两因子均是$K$的多项式，反向乘积相同，故为双侧逆。
\textbf{方法二：}解$x+u(v\T x)=b$。左乘$v\T$得$(1+c)v\T x=v\T b$，由$1+c\ne0$，$v\T x=\frac{v\T b}{1+c}$；代回$x=b-\frac{u(v\T b)}{1+c}$，即给候选逆。若$1+c=0$，则$u\ne0$且$(I+uv\T)u=(1+c)u=0$，确实不可逆，说明分母条件必要。}
{多项式法适合证明，标量消元法说明低秩更新为何便宜。}
{分母为零时行列式为零，不能求逆；$c=0$而$K\ne0$时公式仍成立。}
\begin{warning}
矩阵“等价”“相似”“合同”分别为$PAQ,P\inv AP,P\T AP$，不变量不同。伴随存在不代表逆存在。$A\T A$的秩等于$A$依赖实数域，复数域应换共轭转置。矩阵方程最后用代回检验乘法顺序。
\end{warning}
\section*{章末分层习题}\addcontentsline{toc}{section}{章末分层习题}
\exercise{2-1}{基础}{求$A=\begin{pmatrix}1&2\\0&1\end{pmatrix}$的逆，并计算$A^3$。}
\exercise{2-2}{基础}{设$A,B$均为$n$阶方阵且$AB=I_n$，证明$BA=I_n$。}
\exercise{2-3}{基础}{求$\begin{pmatrix}1&2&3\\2&4&6\end{pmatrix}$的秩。}
\exercise{2-4}{考研风格}{已知$A^2-3A+2I=0$，求$A\inv$。}
\exercise{2-5}{考研风格}{三阶$A$的秩为2，求$\rank\adj A$与$\det\adj A$。}
\exercise{2-6}{考研风格}{解$\diag(2,3)X=\begin{pmatrix}2&4\\3&6\end{pmatrix}$。}
\exercise{2-7}{竞赛拓展}{若$N^3=0$，求$(2I-N)\inv$并验证。}
\exercise{2-8}{竞赛拓展}{设$A$可逆，推导$A+\vu\vv\T$的逆及存在条件。}
\exercise{2-9}{竞赛拓展}{证明任意秩$r$的矩阵都存在秩分解$A=BC$，其中$B$满列秩、$C$满行秩。}
\section*{章末方法速查}\addcontentsline{toc}{section}{章末方法速查}
\begin{quicktable}
伴随求逆&小阶可逆矩阵&二阶、稀疏余子式&伴随要转置\\
增广消元&一般数字矩阵&多个右端&行变换同步\\
多项式求逆&已知低次关系&幂零、幂等&常数项必须可解\\
秩夹逼&参数与抽象矩阵&非零子式、相关式&行列式零不定秩\\
分块逆&可逆主块与Schur补&自然块结构&两个可逆条件\\
分列解方程&系数可奇异&$AX=C$&检查每列相容\\
\end{quicktable}
